% !PRINTED 692
\weberpage{692}
%<<CONT
so ergiebt sich, wenn $\vartheta$ in den Grenzen
\[
0 \leqq \vartheta < 2\pi
\]
genommen wird, bei feststehendem $z$ zu jedem Werthe von $\vartheta$ ein zugehöriges Werthepaar $u_1, u_2$. Ein solcher Winkel $\vartheta$ gehört zu jedem imaginären Paare.

Es gehört dann zu jedem Punkte des Gebietes $S$ ein und nur ein Werthsystem der Variablen $\xi_1, \xi_2, \ldots, \xi_\nu, \vartheta \ldots$ in der Begrenzung

\weq{7}{\begin{aligned}
&0 \leqq \xi_1 < 1, \ldots, 0 \leqq \xi_{\nu-1} < 1 \\
&\xi_\nu < 0, \qquad 0 \leqq \vartheta < 2\pi, \ldots
\end{aligned}}
% [source notation] Both first-row less-than-or-equal signs place the equality bars above the angle; other occurrences on this page place them below. Both variants express the same relation and are represented by \leqq. No damaged-type or reversed-inequality claim is made.

Wenn umgekehrt irgend ein Werthsystem der Variablen $\xi, \vartheta$ in den Grenzen (7) gegeben ist, so erhält man hieraus die entsprechenden Bestandtheile $u_1, u_2$ eines imaginären $y$ eindeutig, von den reellen $y$ aber nur die absoluten Werthe, und es kann also jedem der reellen $y$ noch jedes der beiden Vorzeichen gegeben werden. Die Anzahl dieser möglichen Bestimmungen ist $2^{2\nu-n}$. Ist eine dieser Bestimmungen herausgegriffen, so sind die Variablen $x$ eindeutig bestimmt und sind für jedes endliche Werthsystem der $\xi, \vartheta$ endlich. Für ein gegen $-\infty$ abnehmendes $\xi_\nu$ werden die Werthe der $y$ zum Theil gegen Null convergiren.

Das Gebiet $S$ zerfällt demnach in $2^{2\nu-n}$ Theilgebiete $S_1, S_2, \ldots$, deren jedes durch eine bestimmte Vorzeichen-Combination der reellen $y_1, y_2, \ldots$ charakterisirt ist, und demnach zerfällt das Integral $U$ in ebenso viele Theilintegrale $U_1, U_2, \ldots$

\[
U = U_1+U_2+\cdots
\]

Bei der Ausführung der Integration in einem dieser Bestandtheile, die sich alle als von gleichem Werthe ergeben, fangen wir mit einem imaginären Paare an, wenn ein solches vorhanden ist.

Betrachten wir $u_1, u_2$ als rechtwinkelige Coordinaten in einer Ebene, so können wir $e^{1/2\,z}, \vartheta$ nach (5) als Polarcoordinaten auffassen, und erhalten

\[
\int\int du_1\,du_2 = \frac12\int\int e^z\,dz\,d\vartheta,
\]

worin die Integration in Bezug auf $\vartheta$, die sich von $0$ bis $2\pi$ erstreckt, ausgeführt werden kann. Man findet so

\[
\int\int du_1\,du_2 = \pi\int e^z\,dz.
\]
